\documentclass[a4paper,10pt]{article}
\usepackage[margin=1.5cm]{geometry}
\usepackage{longtable}
\usepackage{listings}
\usepackage[T1]{fontenc}
\usepackage[utf8]{inputenc}

\title{Docker-Compose}
\date{}

\begin{document}
\maketitle


\section*{0. Introducción a \textit{DOCKER-COMPOSE}?}
Es una herramienta de \textbf{Orquestasión de contenedores}, principalmente orinetada en
desarrollo (dev). Aunque su uso no se limita a sólo entorno de desarrollo, hay alternativas más
robustas (como \textbf{KUBERNETES}) para entornos productivos. 

\subsection*{0.1. Instalación (Distribuciones BEBIAN-like )}
Para instalar Docker-Compose en distribuciones DEBIAN, ejecutar los siguientes comandos

\begin{center}
    \begin{lstlisting}
        sudo apt-get update
        sudo apt-get install docker-ce docker-ce-cli containerd.io docker-buildx-plugin docker-compose-plugin   
    \end{lstlisting}
\end{center}

Verifica con

\begin{center}
    \begin{lstlisting}
        docker compose version
    \end{lstlisting}    
\end{center}


\section*{1. El archivo docker-compose.yml}

Este archivo sirve para construir, iniciar y establecer el entorno de los contenedores.
Su sintaxis básica es:


\section*{1. Comandos}

\subsection{}



\end{document}